%%%%%%%%%%  scan idx 94  |  folio 83  %%%%%%%%%%

Sea $A$ una $R$-álgebra. Un $A$-módulo $M$ define un haz aditivo
\[
 \operatorname{Der}_R(-,M)
 \in\operatorname{Ob}(\widetilde{\underline C}_{A|R}),
 \qquad X\longmapsto\operatorname{Der}_R(X,M).
\]
Hacemos la definición siguiente.

\underline{Definición 3.10.} Sea $A$ una $R$-álgebra. Para todo
$A$-módulo $M$, pongamos
\[
 D^i(A|R,M)=H^i(A|R,\operatorname{Der}_R(-,M)).
\]

Calculemos ahora los grupos de grado bajo, a saber, $D^i(A|R,M)$ para
$i=1,2$. Elijamos una flecha sobreyectiva
$P\longrightarrow A$ en $\underline C_{A|R}$, con $P$ un álgebra de
polinomios sobre $R$, y escribamos
\[
 0\longrightarrow I\longrightarrow P\longrightarrow A
 \longrightarrow0.
\]

\underline{Lema 3.10.}
\begin{align*}
 D^1(A|R,M)
 &=\operatorname{Coker}\bigl(
   \operatorname{Der}_R(P,M)\longrightarrow
   \operatorname{Hom}_P(I,M)\bigr),\\
 D^2(A|R,M)
 &=H^1\bigl(A|R,
   \underline H^1_{A|R}(\operatorname{Der}_R(-,M))\bigr).
\end{align*}

\underline{Demostración.} Consideremos el álgebra simplicial
\[
 (*)\qquad
 \cdots\ P_A^{(2)}
 \mathrel{\substack{\longrightarrow\\[-.6ex]\longrightarrow\\[-.6ex]\longrightarrow}}
 P_A^{(1)}
 \mathrel{\substack{\longrightarrow\\[-.6ex]\longrightarrow}}
 P_A^{(0)},
\]
donde
\[
 P_A^{(n)}=
 \overbrace{P\times_A\cdots\times_AP}^{n+1}.
\]
Para todo $n$ hay una sucesión exacta
\[
 0\longrightarrow I^{(n)}\longrightarrow P_A^{(n)}
 \longrightarrow A\longrightarrow0,
 \qquad
 I^{(n)}=\overbrace{I\times\cdots\times I}^{n+1}.
\]
Pongamos $\overline I=I/I^2$, y denotemos por $\Delta$ el morfismo
diagonal. Se comprueba inmediatamente que, para todo $n$, la sucesión
\[
 0\longrightarrow
 \operatorname{Hom}_A
 \bigl(\overline I^{(n)}/\Delta(\overline I),M\bigr)
 \longrightarrow\operatorname{Der}_R(P_A^{(n)},M)
 \xrightarrow{\operatorname{Der}(\Delta,M)}
 \operatorname{Der}_R(P,M)\longrightarrow0
 \tag{a}
\]
es exacta. También tenemos la sucesión exacta
\[
 0\longrightarrow
 \operatorname{Hom}_A
 \bigl(\overline I^{(n)}/\Delta(\overline I),M\bigr)
 \longrightarrow\operatorname{Hom}_A(\overline I^{(n)},M)
 \xrightarrow{\operatorname{Hom}(\Delta,M)}
 \operatorname{Hom}_A(\overline I,M)\longrightarrow0.
 \tag{b}
\]
Las sucesiones (a) y (b) dan lugar a las correspondientes sucesiones
exactas de $A$-módulos simpliciales. Puesto que
$\{I\longrightarrow0\}$ y $\{I\longrightarrow I\}$ admiten secciones,
\[
 H^i(\{I\longrightarrow0\},\operatorname{Hom}_A(-,M))
 =H^i(\{I\longrightarrow I\},\operatorname{Hom}_A(-,M))=0
 \qquad(i>0).
\]

%%%%%%%%%%  scan idx 95  |  folio 84  %%%%%%%%%%

Se sigue que
\[
\left\{
\begin{aligned}
0\longrightarrow\operatorname{Der}_R(A,M)
&\longrightarrow\operatorname{Der}_R(P,M)
\longrightarrow\operatorname{Hom}_A(\overline I,M)\\
&\longrightarrow
H^1(\{P\longrightarrow A\},\operatorname{Der}_R(-,M))
\longrightarrow0
&&\text{es exacta},\\
H^i(\{P\longrightarrow A\},\operatorname{Der}_R(-,M))
&=0 &&(i>1).
\end{aligned}
\right.
\]
Por tanto, según 3.1(b),
\begin{align*}
D^1(A|R,M)
 &=\check H^1(A|R,\operatorname{Der}_R(-,M))\\
 &=H^1(\{P\longrightarrow A\},\operatorname{Der}_R(-,M))\\
 &=\operatorname{Coker}\bigl(
   \operatorname{Der}_R(P,M)\longrightarrow
   \operatorname{Hom}_A(\overline I,M)\bigr)\\
 &=\operatorname{Coker}\bigl(
   \operatorname{Der}_R(P,M)\longrightarrow
   \operatorname{Hom}_P(I,M)\bigr),\\
D^2(A|R,M)
 &=H^1\bigl(A|R,
   \underline H^1_{A|R}(\operatorname{Der}_R(-,M))\bigr).
\end{align*}
\hfill$\square$

Elijamos un $P$-módulo libre $F$ y un morfismo $P$-lineal
$\gamma:F\longrightarrow P$ tal que $\gamma(F)=I$. A $\gamma$ se
asocia un complejo de Koszul que, por abuso de notación, denotaremos
por $K_I$.
\par
